% Generated by make_inputs_table.py -- DO NOT EDIT BY HAND.
% Rebuild:  PYTHONPATH=. python3 make_inputs_table.py
%           pdflatex inputs_table && bibtex inputs_table && pdflatex inputs_table x2
\documentclass[11pt,a4paper]{article}
\usepackage[margin=2.2cm]{geometry}
\usepackage{amsmath,amssymb,longtable,xcolor}
\usepackage[colorlinks=true,linkcolor=black,citecolor=blue,urlcolor=blue]{hyperref}
\title{Inventory of inputs, with sources}
\author{Compiled for the \texttt{ionization\_yield} / \texttt{photon\_vs\_electron} project}
\date{2026-09-14}
\begin{document}\maketitle

\noindent
Every number this project takes from outside itself, with its symbol, value,
range, units and source. \textbf{Values are read from the code that uses them}
by \texttt{make\_inputs\_table.py}, so the table cannot drift from the
calculation; only the symbols, units and citations are written by hand.

\paragraph{What the status column is for.} It separates four very different
kinds of number that a bare citation would make look alike: a value checked
numerically against the paper it came from; a value whose paper we have never
opened; a value you chose; and a value that is conventional in the field but
for which \emph{this project has established no source at all}. The last
category is the one worth acting on.

\section{Inventory of inputs}\label{sec:inputs}
Every externally supplied number in this project. Values are read from the code that uses them, so this table cannot drift from the calculation. Status: \textbf{V} verified numerically against the source, which is in \texttt{papers/}; C cited but the source is not in the tree and the implementation is unchecked; D derived here; S scanned; U your stated choice; \textcolor{red}{\textbf{X}} conventional value whose source this project has \emph{not} established.
\begin{longtable}{llllll}
\hline
symbol & value & range & units & source & st. \\
\hline\endhead
\multicolumn{6}{l}{\textbf{Cosmology}}\\[2pt]
$h$ & 0.6766 & -- & -- & \cite{Planck2020} & C \\
$\Omega_m$ & 0.3111 & -- & -- & \cite{Planck2020} & C \\
$\Omega_\Lambda$ & 0.6889 & -- & -- & \cite{Planck2020} & D \\
\multicolumn{6}{p{0.95\textwidth}}{\footnotesize\quad flat by assumption}\\[2pt]
$\Omega_b h^2$ & 0.02242 & -- & -- & \cite{Planck2020} & C \\
$\rho_{\rm crit}/h^2$ & 1.87834e-29 & -- & g\,cm$^{-3}$ & \cite{PDG2022} & C \\
$T_{\rm CMB,0}$ & 2.7255 & -- & K & \cite{Fixsen2009} & C \\
$Y_P$ & 0.245 & -- & -- & \cite{Planck2020} & C \\
\multicolumn{6}{p{0.95\textwidth}}{\footnotesize\quad BBN + Planck}\\[2pt]
$X_H$ & 0.755 & -- & -- & \cite{Planck2020} & D \\
\multicolumn{6}{p{0.95\textwidth}}{\footnotesize\quad $1-Y_P$}\\[2pt]
$H(z)$ & Planck18 & -- & s$^{-1}$ & \cite{Planck2020} & C \\
\multicolumn{6}{p{0.95\textwidth}}{\footnotesize\quad \texttt{astropy}; includes radiation}\\[2pt]
\multicolumn{6}{l}{\textbf{IGM target}}\\[2pt]
$n_{\rm H}(z{=}10)$ & 2.5287e-04 & -- & cm$^{-3}$ & \cite{Planck2020} & D \\
\multicolumn{6}{p{0.95\textwidth}}{\footnotesize\quad proper; from $\Omega_b h^2$, $X_H$}\\[2pt]
$N_{\rm HI,norm}$ & 1759.467095 & -- & m$^{-3}$ & \cite{Planck2020} & D \\
\multicolumn{6}{p{0.95\textwidth}}{\footnotesize\quad $n_{\rm HI}(z)=N((1+z)/21)^3$; harmonised}\\[2pt]
$x_e$ & 1e-04 & -- & -- & --- & U \\
\multicolumn{6}{p{0.95\textwidth}}{\footnotesize\quad held static; your value}\\[2pt]
$B_0$ & 1.0e-13 & -- & T & --- & U \\
\multicolumn{6}{p{0.95\textwidth}}{\footnotesize\quad 1 nG comoving; $B\propto(1+z)^2$}\\[2pt]
$Z$ & 1 & -- & -- & --- & U \\
\multicolumn{6}{p{0.95\textwidth}}{\footnotesize\quad pure hydrogen target}\\[2pt]
$\delta_{\rm Gould}$ & 0.5 & -- & -- & \cite{Gould1972} & C \\
\multicolumn{6}{p{0.95\textwidth}}{\footnotesize\quad max fractional transfer}\\[2pt]
$f_{\rm therm}$ & 1.5 & -- & -- & --- & U \\
\multicolumn{6}{p{0.95\textwidth}}{\footnotesize\quad $K_{\rm term}=f\,kT_{\rm CMB}$}\\[2pt]
$n_{\rm max}^{\rm exc}$ & 10 & -- & -- & \cite{Stone2002} & C \\
\multicolumn{6}{p{0.95\textwidth}}{\footnotesize\quad levels $n=2\ldots n_{\rm max}$}\\[2pt]
\multicolumn{6}{l}{\textbf{Atomic physics}}\\[2pt]
$E_{\rm th}$(HI) & 13.598434600 & -- & eV & \cite{NIST_ASD} & \textbf{V} \\
$R_\infty$ & 13.605693123 & -- & eV & \cite{Tiesinga2021} & \textbf{V} \\
$B$ (Kim Table I) & 13.6057 & -- & eV & \cite{KimRudd1994} & \textbf{V} \\
\multicolumn{6}{p{0.95\textwidth}}{\footnotesize\quad the Rydberg; tension with NIST at 0.05\%}\\[2pt]
$U/B$ & 1 & -- & -- & \cite{KimRudd1994} & \textbf{V} \\
\multicolumn{6}{p{0.95\textwidth}}{\footnotesize\quad virial, H(1s)}\\[2pt]
$N_i$ & 0.434306 & -- & -- & \cite{KimRudd1994} & \textbf{V} \\
\multicolumn{6}{p{0.95\textwidth}}{\footnotesize\quad continuum oscillator strength}\\[2pt]
$M^2$ & 0.2834 & -- & -- & \cite{KimRudd1994} & \textbf{V} \\
$a_i$ (d$f$/d$w$) & 4 coeffs & -- & -- & \cite{KimRudd1994} & \textbf{V} \\
\multicolumn{6}{p{0.95\textwidth}}{\footnotesize\quad Table I, H(1s)}\\[2pt]
$\sigma_{\rm ion}$ & RBED eq.~(20) & -- & m$^2$ & \cite{Kim2000} & \textbf{V} \\
\multicolumn{6}{p{0.95\textwidth}}{\footnotesize\quad verified to $1.4\times10^{-9}$}\\[2pt]
d$\sigma$/d$w$ & BED, non-rel. & -- & m$^2$eV$^{-1}$ & \cite{KimRudd1994} & \textbf{V} \\
\multicolumn{6}{p{0.95\textwidth}}{\footnotesize\quad relativistic factor omitted: +25\% at 1 MeV}\\[2pt]
$\sigma_{\rm exc}$ & Stone \& Kim & -- & m$^2$ & \cite{Stone2002} & C \\
Coulomb loss & Gould (1972) & -- & eV\,s$^{-1}$ & \cite{Gould1972} & C \\
IC / brems & BG70 & -- & eV\,s$^{-1}$ & \cite{Blumenthal1970} & C \\
\multicolumn{6}{p{0.95\textwidth}}{\footnotesize\quad eq.~(2.42) spectrum verified; loss rates not}\\[2pt]
$I_{\rm exc}$ & 14.99 & -- & eV & \cite{ICRU1984} & C \\
\multicolumn{6}{p{0.95\textwidth}}{\footnotesize\quad mean excitation energy, route 1 only}\\[2pt]
$\alpha_B(2\times10^4\,$K$)$ & 1.430e-13 & -- & cm$^3$s$^{-1}$ & \cite{OsterbrockFerland2006} & \textbf{V} \\
\multicolumn{6}{p{0.95\textwidth}}{\footnotesize\quad confirmed vs \cite{HuiGnedin1997} to 0.16\%}\\[2pt]
$\chi$ & 1.08 & -- & -- & \cite{MadauDickinson2014} & \textbf{V} \\
\multicolumn{6}{p{0.95\textwidth}}{\footnotesize\quad derived here as $1+(Y_P/4)/X_H=1.0811$}\\[2pt]
\multicolumn{6}{l}{\textbf{Stellar source}}\\[2pt]
$\log_{10}\rho_{\rm UV}$ & 25.12 & $+0.07/-0.08$ & erg\,s$^{-1}$Hz$^{-1}$cMpc$^{-3}$ & \cite{Donnan2024} & \textbf{V} \\
\multicolumn{6}{p{0.95\textwidth}}{\footnotesize\quad Table 3, $z=10$ row}\\[2pt]
$\log_{10}\xi_{\rm ion}$ & 25.28 & $\pm0.42$ dex & Hz\,erg$^{-1}$ & \cite{Llerena2025} & \textbf{V} \\
\multicolumn{6}{p{0.95\textwidth}}{\footnotesize\quad median at $z\simeq7.14$; scatter is observed, not error}\\[2pt]
$K_{\rm UV}$ & 1.150e-28 & -- & $M_\odot$yr$^{-1}$/(erg\,s$^{-1}$Hz$^{-1}$) & \cite{MadauDickinson2014} & \textbf{V} \\
\multicolumn{6}{p{0.95\textwidth}}{\footnotesize\quad eq.~(10); Salpeter $0.1$--$100\,M_\odot$}\\[2pt]
$\alpha_{\rm IMF}$ & 2.35 & -- & -- & \cite{Salpeter1955} & \textbf{V} \\
\multicolumn{6}{p{0.95\textwidth}}{\footnotesize\quad fitted for $0.4$--$10\,M_\odot$ only}\\[2pt]
$m_{\min},m_{\max}$ & 0.1, 100.0 & -- & $M_\odot$ & \cite{MadauDickinson2014} & C \\
\multicolumn{6}{p{0.95\textwidth}}{\footnotesize\quad range adopted to match $K_{\rm UV}$}\\[2pt]
$m_{\rm SN}$ & 8.0 & -- & $M_\odot$ & --- & U \\
\multicolumn{6}{p{0.95\textwidth}}{\footnotesize\quad core-collapse threshold}\\[2pt]
SN per $M_\odot$ & 7.4218e-03 & -- & $M_\odot^{-1}$ & \cite{Salpeter1955} & D \\
\multicolumn{6}{p{0.95\textwidth}}{\footnotesize\quad IMF integral, this work}\\[2pt]
$\alpha_{\rm SED}$ & 2 & 1--3 & -- & \cite{MarquesChaves2026} & C \\
\multicolumn{6}{p{0.95\textwidth}}{\footnotesize\quad $f_\nu\propto\nu^{-\alpha}$; BPASS spectrum unobtainable}\\[2pt]
LyC band & 1--4 & -- & Ryd & --- & U \\
\multicolumn{6}{p{0.95\textwidth}}{\footnotesize\quad conventional}\\[2pt]
$T_{\rm eff}$ & 5e+04 & -- & K & --- & U \\
\multicolumn{6}{p{0.95\textwidth}}{\footnotesize\quad blackbody, comparison case only}\\[2pt]
\multicolumn{6}{l}{\textbf{Cosmic-ray channel}}\\[2pt]
$E_{\rm SN}$ & 1e+51 & -- & erg & --- & \textcolor{red}{\textbf{X}} \\
\multicolumn{6}{p{0.95\textwidth}}{\footnotesize\quad standard core-collapse budget; source not established here}\\[2pt]
$\epsilon_{\rm CR}$ & 0.1 & $0.1$--$0.2$ & -- & \cite{CaprioliSpitkovsky2014} & C \\
\multicolumn{6}{p{0.95\textwidth}}{\footnotesize\quad hybrid simulations: $10$--$20\%$ of bulk kinetic energy into non-thermal protons at parallel/quasi-parallel strong shocks, dropping to ZERO when quasi-perpendicular; $\simeq10\%$ inferred for Tycho. An obliquity-AVERAGED value would be lower, so $0.1$ favours the CR channel}\\[2pt]
$f_e=K_{ep}$ & 0.01 & $1$--$3\times10^{-3}$ (source) & -- & \cite{Park2015} & C \\
\multicolumn{6}{p{0.95\textwidth}}{\footnotesize\quad \textbf{the project uses the value measured at Earth, not the value the source derives}: \cite{Park2015} give $K_{ep}\simeq1$--$3\times10^{-3}$ at realistic SNR shock speeds and $1.6\times10^{-3}$ for Tycho, so $0.01$ is $3$--$10\times$ high and favours the CR channel. $K_{ep}$ is defined at fixed MOMENTUM and is independent of $p$}\\[2pt]
$p$ & 2.2 & -- & -- & --- & U \\
\multicolumn{6}{p{0.95\textwidth}}{\footnotesize\quad your choice; DSA strong-shock predicts 2}\\[2pt]
$E_{\min},E_{\max}$ & 1e+03, 1e+12 & -- & eV & --- & U \\
\multicolumn{6}{p{0.95\textwidth}}{\footnotesize\quad your choice; spectrum cancels from $\zeta_e$}\\[2pt]
\multicolumn{6}{l}{\textbf{Reionization budget}}\\[2pt]
$C_{\rm IGM}(z)$ & $1+43z^{-1.71}$ & -- & -- & \cite{Pawlik2009} & \textbf{V} \\
\multicolumn{6}{p{0.95\textwidth}}{\footnotesize\quad eq.~(A1), $C_{100}$, r19.5 run; $43=e^{3.76}$; fit valid $6\le z\le20$}\\[2pt]
$\langle t_{\rm rec}\rangle$ & eq.~(24) & -- & Gyr & \cite{MadauDickinson2014} & \textbf{V} \\
\multicolumn{6}{p{0.95\textwidth}}{\footnotesize\quad 3.2 Gyr normalisation reproduced to 1.6\%}\\[2pt]
\multicolumn{6}{l}{\textbf{Scanned ranges}}\\[2pt]
$f_{\rm esc}$ & 0.1 & $10^{-6}$--$1$ & -- & --- & S \\
\multicolumn{6}{p{0.95\textwidth}}{\footnotesize\quad canonical value is a convention, not a measurement}\\[2pt]
$\epsilon_{\rm CR}f_e$ & 1e-03 & $10^{-6}$--$1$ & -- & --- & S \\
\multicolumn{6}{p{0.95\textwidth}}{\footnotesize\quad bounded above by unity on energy grounds}\\[2pt]
$\rho_{\rm SFR}$ scale & 1 & $10^{-3}$--$1$ & -- & --- & S \\
\multicolumn{6}{p{0.95\textwidth}}{\footnotesize\quad ignorance band at $z=20$ only, where $\rho_{\rm UV}$ is unmeasured}\\[2pt]
\multicolumn{6}{l}{\textbf{Source spectra}}\\[2pt]
$\alpha_{\rm EUV}$ (RQ) & 1.57 & $\pm0.17$ & -- & \cite{Telfer2002} & \textbf{V} \\
\multicolumn{6}{p{0.95\textwidth}}{\footnotesize\quad HST composite, rest-frame $500$--$1200$\,\AA}\\[2pt]
$\alpha_{\rm EUV}$ (RL) & 1.96 & $\pm0.12$ & -- & \cite{Telfer2002} & \textbf{V} \\
\multicolumn{6}{p{0.95\textwidth}}{\footnotesize\quad radio-loud quasars; full composite $1.76\pm0.12$}\\[2pt]
$\alpha_{\rm AGN}$ & 1.5 & -- & -- & \cite{Graziani2018} & \textbf{V} \\
\multicolumn{6}{p{0.95\textwidth}}{\footnotesize\quad EoR quasar template; corroborates \cite{Telfer2002}}\\[2pt]
AGN band top & 3 & -- & keV & \cite{Graziani2018} & \textbf{V} \\
\multicolumn{6}{p{0.95\textwidth}}{\footnotesize\quad $13.6$\,eV--$3$\,keV; above 4 Ryd is an extrapolation}\\[2pt]
$\Gamma_1$ (ULX) & 2.1 & 1.38--3.1 & -- & \cite{Gladstone2009} & \textbf{V} \\
\multicolumn{6}{p{0.95\textwidth}}{\footnotesize\quad Table~6; single-PO fits $1.6<\Gamma<3.3$}\\[2pt]
$E_{\rm break}$ (ULX) & 5.0 & 3.5--7 & keV & \cite{Gladstone2009} & \textbf{V} \\
\multicolumn{6}{p{0.95\textwidth}}{\footnotesize\quad broken PL preferred $>98\%$ in 11 of 12; fits are $2$--$10$\,keV only}\\[2pt]
$\Delta\Gamma$ (ULX) & 1.5 & 1--2 & -- & \cite{Gladstone2009} & \textbf{V} \\
\multicolumn{6}{p{0.95\textwidth}}{\footnotesize\quad steepening above the break}\\[2pt]
$kT_e,\tau$ (ULX corona) & 1--3, 6--80 & -- & keV, -- & \cite{Gladstone2009} & \textbf{V} \\
\multicolumn{6}{p{0.95\textwidth}}{\footnotesize\quad cool, optically THICK; global $\chi^2$ min in 10/12}\\[2pt]
$\Gamma$ (BHB hard) & 1.7 & $<2.1$ & -- & \cite{Gladstone2009} & \textcolor{red}{\textbf{X}} \\
\multicolumn{6}{p{0.95\textwidth}}{\footnotesize\quad \textbf{the typical $\Gamma\simeq1.7$ is second-hand}; only $\Gamma<2.1$ for the low/hard state is first-hand (\S4.1)}\\[2pt]
$p$ per class & 2.2 & -- & -- & --- & U \\
\multicolumn{6}{p{0.95\textwidth}}{\footnotesize\quad held at the project fiducial: no citable per-class CR index at $z\sim10$}\\[2pt]
\multicolumn{6}{l}{\textbf{Source luminosities}}\\[2pt]
AGN LyC fraction & 0.17--0.33 & -- & -- & \cite{Asthana2024} & \textbf{V} \\
\multicolumn{6}{p{0.95\textwidth}}{\footnotesize\quad $17\%$ (\cite{Asthana2024}) to $\le1/3$ (\cite{Jiang2025}), both at $z\sim7.5$; carried to $z=10$, conservative for photons}\\[2pt]
$L_{2-10}/$SFR & 1.8e+39 & 1.4e+39--2.3e+39 & erg\,s$^{-1}$/($M_\odot$yr$^{-1}$) & \cite{Lehmer2016} & \textbf{V} \\
\multicolumn{6}{p{0.95\textwidth}}{\footnotesize\quad eq.~(6), $\times(1+z)^{1}$; exponent fixed at 1}\\[2pt]
$L_{\rm mech}$ (microquasar) & 1e+39--1e+40 & -- & erg\,s$^{-1}$ & \cite{Mirabel2011} & \textbf{V} \\
\multicolumn{6}{p{0.95\textwidth}}{\footnotesize\quad SS\,433 and the S26 microquasar; TOTAL injected, not the electron share}\\[2pt]
$E_{\rm mech}$ per object & 1e+54 & -- & erg & \cite{Mirabel2011} & \textbf{V} \\
\multicolumn{6}{p{0.95\textwidth}}{\footnotesize\quad whole lifetime; $\gg$ a core-collapse SN}\\[2pt]
$M_{\rm UV}$ (U37126) & -20.1 & $\pm0.05$ & AB & \cite{MarquesChaves2026} & \textbf{V} \\
\multicolumn{6}{p{0.95\textwidth}}{\footnotesize\quad $z=10.255$, lensed $\mu\simeq2.2$; the per-object stellar anchor}\\[2pt]
$f_{\rm jet,e}$ & 1e-03--1 & scanned & -- & --- & S \\
\multicolumn{6}{p{0.95\textwidth}}{\footnotesize\quad jet efficiency $\times$ leptonic fraction, scanned: the proton content of jets is not observationally constrained}\\[2pt]
$f_{\rm dep}$ & 1 & -- & -- & --- & U \\
\multicolumn{6}{p{0.95\textwidth}}{\footnotesize\quad assumption A11: every injected electron deposits --- a CEILING}\\[2pt]
\multicolumn{6}{l}{\textbf{Derived (this work)}}\\[2pt]
$W_e$ & 48.038 & -- & eV\,ion$^{-1}$ & --- & D \\
\multicolumn{6}{p{0.95\textwidth}}{\footnotesize\quad $\langle E\rangle/\langle N\rangle$ of the fiducial CR spectrum}\\[2pt]
$W_\gamma$ (stellar) & 21.07 & -- & eV\,ion$^{-1}$ & --- & D \\
\multicolumn{6}{p{0.95\textwidth}}{\footnotesize\quad $W_\gamma/W_e=0.439$; IGM transparency applied}\\[2pt]
$W_\gamma$ (AGN) & 27.41 & -- & eV\,ion$^{-1}$ & --- & D \\
\multicolumn{6}{p{0.95\textwidth}}{\footnotesize\quad $W_\gamma/W_e=0.571$; IGM transparency applied}\\[2pt]
$W_\gamma$ (ULX) & 38.78 & -- & eV\,ion$^{-1}$ & --- & D \\
\multicolumn{6}{p{0.95\textwidth}}{\footnotesize\quad $W_\gamma/W_e=0.807$; IGM transparency applied}\\[2pt]
$W_\gamma$ (BHB hard) & 146.83 & -- & eV\,ion$^{-1}$ & --- & D \\
\multicolumn{6}{p{0.95\textwidth}}{\footnotesize\quad $W_\gamma/W_e=3.057$; IGM transparency applied}\\[2pt]
$\tau_{\rm IGM}$(1 keV) & 1.93 & -- & -- & --- & D \\
\multicolumn{6}{p{0.95\textwidth}}{\footnotesize\quad over one Hubble length at $z=10$; thick below $\sim500$\,eV, transparent above $\sim3$\,keV}\\[2pt]
$\tau_{\rm IGM}$(3 keV) & 0.0476 & -- & -- & --- & D \\
\multicolumn{6}{p{0.95\textwidth}}{\footnotesize\quad over one Hubble length at $z=10$; thick below $\sim500$\,eV, transparent above $\sim3$\,keV}\\[2pt]
$\tau_{\rm IGM}$(10 keV) & 0.00077 & -- & -- & --- & D \\
\multicolumn{6}{p{0.95\textwidth}}{\footnotesize\quad over one Hubble length at $z=10$; thick below $\sim500$\,eV, transparent above $\sim3$\,keV}\\[2pt]
$t_H(z{=}10)$ & 0.7086 & -- & Gyr & \cite{Planck2020} & D \\
\multicolumn{6}{p{0.95\textwidth}}{\footnotesize\quad astropy Planck18, complete $H(z)$}\\[2pt]
\hline
\end{longtable}

\section{The fiducial top-heavy IMF}\label{sec:topheavy}
For the scenario considered here --- Population~II stars in galaxies at
$z=10$ --- the fiducial top-heavy IMF adopted is that of
Jeong, Jeon, Song \& Bromm~\cite{Jeong2025}, their eq.~(2):
\begin{equation}
  \xi(m) = \frac{\mathrm{d}N}{\mathrm{d}m} = \xi_0\,m^{-\alpha},
  \qquad \alpha = 1.3,
  \qquad 1\,M_\odot \le m \le 100\,M_\odot .
  \label{eq:topheavy}
\end{equation}
It is chosen over the alternatives for four reasons. \emph{First}, it is
stated explicitly as a top-heavy IMF \emph{for Pop~II stars}, not for
Pop~III, which is the population this project's galaxies contain.
\emph{Second}, the galaxies it is applied to form at $z=10$, the redshift of
this calculation. \emph{Third}, it is the same functional form as the
Salpeter law already in use, $\mathrm{d}N/\mathrm{d}m \propto m^{-\alpha}$,
so it substitutes directly into the supernova-rate integral with no change of
machinery --- only $\alpha$ and $m_{\min}$ move. \emph{Fourth}, the physical
motivation is the one relevant here: a raised temperature floor from
low-metallicity gas radiatively coupled to a CMB at $T=2.73(1+z)$~K, which
lifts the Jeans mass and hence the characteristic stellar
mass~\cite{Larson1998}. Their default Pop~II IMF, for comparison, is
Chabrier~\cite{Chabrier2003} over $0.1$--$100\,M_\odot$.

\paragraph{Not yet adopted.} Equation~\eqref{eq:topheavy} is recorded here as
the fiducial to test against; it has \emph{not} been substituted into any
result in this project. Doing so would move three of the entries in
\S\ref{sec:inputs} at once --- $\alpha_{\rm IMF}$, the mass range, and the
supernova rate per $M_\odot$ --- and would also invalidate $K_{\rm UV}$,
which \cite{MadauDickinson2014} quote for a Salpeter IMF over $0.1$--$100\,M_\odot$.
Those three quantities are not independent, and that coupling is the open
question this inventory was built to expose.

\bibliographystyle{plain}
\bibliography{../papers/references}
\end{document}

